\subsection{Syntax}
\label{sec:syntax}

The Pat syntax we are working with is close to the original definitions by
Fowler et al.\ \cite{fowlerSpecialDeliveryProgramming2023}, making use of
de Bruijn indices \cite{debruijnLambdaCalculusNotation1972} to represent
variables.
We use de Bruijn indices as they are well-studied for proof assistants, and in
particular Rocq \cite{ambalHOPiCoq2021,starkMechanisingSyntaxBinders2019}.
We make use of the Rocq library \emph{dblib} \cite{pottierDblib2021} that
automatically handles de Bruijn indices and operations on them such as shifting.

To ease the burden of proving properties about messages, we restrict messages to
only carry a single payload, and consider only functions that take a single
argument.  We expect that extending our proofs to include arbitrary numbers of
payloads and function arguments will not significantly affect the proofs.

\begin{figure}[t]
  \centering
  \begin{tabular}{lccl}
    % Program definition \; & $\prog$ & & \\
    Definition names (\rocqlink{https://edgardschi.github.io/mailbox-types-rocq/MailboxTypes.Syntax.html\#IDefinitionName}) \; & $\defname$ & & \\
    Function definitions (\rocqlink{https://edgardschi.github.io/mailbox-types-rocq/MailboxTypes.Syntax.html\#FunctionDefinition}) \; & $d$ & $::=$ & $\FunctionDef$\\
    Booleans \; & $b$ & $::=$ & $\mathsf{true} \mid \mathsf{false}$ \\
    Values (\rocqlink{https://edgardschi.github.io/mailbox-types-rocq/MailboxTypes.Syntax.html\#Value}) \; & $v$ & $::=$ & $b \mid \unitvalue \mid \Var{n}$ \\
    Terms (\rocqlink{https://edgardschi.github.io/mailbox-types-rocq/MailboxTypes.Syntax.html\#Term}) \; & $t$ & $::=$ & $v \mid \TLet{t_1}{t_2} \mid \defname(v) \mid \TSpawn{t}$\\
    & & $\mid$ & $\TNew \mid \TSend{v_1}{m}{v_2} \mid \TGuard{v}{E}{\overrightarrow{g}}$\\
    Guards (\rocqlink{https://edgardschi.github.io/mailbox-types-rocq/MailboxTypes.Syntax.html\#Guard}) \; & $g$ & $::=$ & $\GFail \mid \GFree{t} \mid \GReceive{m}{t}$\\
    Environments (\rocqlink{https://edgardschi.github.io/mailbox-types-rocq/MailboxTypes.Environment.html\#Env}) \; & $\Gamma$ & $::=$ & $\emptyset \mid \bot,\Gamma \mid A,\Gamma$
  \end{tabular}

  \caption{The syntax of Pat}
  \label{fig:pat-syntax}
\end{figure}

Figure \ref{fig:pat-syntax} shows our syntax of Pat.
A value $v$ is either a Boolean value, the unit value $\unitvalue$, or a variable $\Var{n}$ represented as a de Bruijn index, with natural number $n$.
Including variables into the definition of values is fairly nonstandard.
The inspiration for this approach is the \emph{fine-grain call-by-value} evaluation strategy \cite{levyModellingEnvironmentsCallbyvalue2003},
that makes a syntactic difference between values and computations.
Values, including variables that can also represent mailbox names, are immutable and cannot be reduced, while
other terms produce values and can be reduced further.

A term is either a value, a let-expression or a function application $\defname(v)$ of function $\defname$ on value $v$.
Additionally, there are terms specifically for interacting with mailboxes and processes.
The term $\TSpawn{t}$ can be used to create a new process running term $t$.
Dynamic mailbox name creation is done with the $\TNew$ keyword, generating a fresh mailbox name.
Term $\TSend{v_1}{\msg{m}}{v_2}$ sends message $\msg{m}$ with payload value $v_2$ to the mailbox with name $v_1$.
A guard term $\TGuard{v}{E}{\overrightarrow{g}}$ asserts that the mailbox associated with name $V$ has pattern $E$, and invokes some guard from the list of guards $\overrightarrow{g}$.
Guards perform actions on a mailbox:
the guard $\GFail$ is triggered when a mailbox receives an unexpected message,
and should be evaluated by a well-typed program.
In contrast, guard $\GFree{t}$ is triggered when a mailbox is empty and not
referenced elsewhere in the system;
the corresponding mailbox is freed and we proceed with continuation $t$.
Finally, the guard $\GReceive{m}{t}$ is triggered when a mailbox contains the message $m$.
The contents of the message are then available in the continuation $t$.

A Pat program $\mathcal{P} = (\mathcal{S}, \mathcal{D}, t_0)$ (\rocqlink{https://edgardschi.github.io/mailbox-types-rocq/MailboxTypes.Syntax.html\#Prog}) is a tuple consisting of a signature $\mathcal{S}$ for mapping messages to their payload type,
a mapping $\mathcal{D}$ from definition names to function definitions, and an initial term $t_0$.
For simplicity, given a program $\mathcal{P}$, we write $\mathcal{S}_\mathcal{P}$ for the signature, and $\mathcal{D}_\mathcal{P}$ for the function definitions.

\paragraph{Mechanization.}
The mechanization of the syntax of Pat is close to its pen-and-paper counterpart,
with the exception of using de Bruijn indices.
Defining the syntax poses no problem after settling on the representation of variables.
While the choice of using de Bruijn indices is standard \cite{ambalHOPiCoq2021,tiroreMultipartyAsynchronousSession2025,starkMechanisingSyntaxBinders2019}, it comes with the downside of extensive reasoning about indices and lifting.
While \emph{dblib} alleviates some of this burden, a substantial amount of manual work is still required.

Definition names are treated similar to message tags, being represented as an arbitrary type equipped decidable equality.
